% Volume IX: Institutional Design
% Part XVII. Designing for Delayed Truth
% Chapter 99: Bounded Transparency
% Spine: proof-spines/ch099-spine.md

\chapter{Bounded Transparency}
\label{ch:ch099}

Transparency should be designed by reference to purpose, audience, timing, and contestability rather than treated as a choice between total secrecy and universal publication. The question is not whether to disclose everything, but how to disclose enough, to the right people, at the right time, and in a form that permits challenge without manufacturing needless exposure.

Between opacity and full publicity lies a whole repertoire of intermediate forms: delayed release, redaction, confidential review, selective disclosure, and independent escrow. These are not evasions of accountability, but design options for making scrutiny compatible with privacy, safety, and procedural fairness. Bounded transparency is therefore the architecture of answerability under conditions where disclosure itself can help or harm.
